\section{Finite signed diffuse data and an exact relative obstacle problem}
\label{sec:r11-diffuse-obstacle-contact}

Keep the original arbitrary symmetric middle-positive separately
closed $(n-1,n-1)$-current $S$, the analytic affinoid neighborhood
$\Omega$, and the ordinary part domain $V\Subset\Omega$ from the
signed diffuse Dirichlet theorem.  Write $(H,\mathcal E,c_V,G,\lambda)$
for its faithful part form and put
\[
                              Az=-D(zS).
 \tag{OB1}
\]
Thus $G\sigma=-u_\sigma$ for a finite signed energy-diffuse measure
$\sigma$, and $A(G\sigma)=\sigma$ on every original test.
All functions below use their canonical energy-quasicontinuous
representatives on $F$.  Inequalities mean quasi-everywhere on $F$;
no pointwise boundary condition or ambient classical psh condition
is substituted for the specified part-form boundary class.

\begin{lemma}[Positive parts are Green potentials]
\label{lem:r11-diffuse-positive-part-Green}
Let $\sigma$ be any finite signed energy-diffuse Radon measure and
let $z=G\sigma$.  Then $z^+$ is the Green potential of its own finite
signed diffuse curvature and
\[
 \begin{split}
 Az^+&=\mathbf1_{\{\widetilde z>0\}}\sigma-L_0(z;S),\\
 z^+&=G(Az^+),\qquad L_0(z;S)\ge0.
 \end{split}
 \tag{OB2}
\]
In particular,
\[
        \mathbf1_{\{\widetilde z>0\}}\sigma\le0
                              \quad\Longrightarrow\quad z\le0.
 \tag{OB3}
\]
These assertions do not require $z$ or $z^+$ to belong to $H$.
\end{lemma}

\noindent\textbf{Status: \texttt{DIRECT\_COROLLARY}} of the full
signed exact-level and global curvature-variation theorems, with the
Green representation proved explicitly below.

\begin{proof}
The signed Kato theorem gives the first line of (OB2), and its level
measure is finite and diffuse by (SD10)--(SD13).  For $M>0$ put
$v_M=(T_Mz)^+\in H$.  The actual scalar identity
$v_M=z^+-(z-M)^+$ and the signed level formula give
\[
 D(v_MS)=\mathbf1_{\{0<\widetilde z\le M\}}D(zS)
                                    +L_0(z;S)-L_M(z;S).
 \tag{OB4}
\]
Consequently
\[
 \|D(v_MS)-D(z^+S)\|_{\rm TV}
                   \le2|\sigma|\{\widetilde z>M\}\longrightarrow0.
\]
Each $v_M$ is an $H$ vector with finite signed Radon curvature, so
Lemma~\ref{lem:r11-bounded-H-tests-diffuse} identifies it with the
transposition solution of that curvature.  Total-variation to
$L^1(\lambda)$ stability then gives
$v_M\to u_{D(z^+S)}$.  Also $v_M\to z^+$ in $L^1(\lambda)$ by
dominated convergence, since $|v_M|\le|z|$.  This proves the second
line of (OB2), including the canonical representative: its bounded
clippings agree as $H$ vectors and hence quasi-everywhere.

Under the hypothesis in (OB3), (OB2) gives $Az^+\le0$.
Order preservation of $G$ therefore gives $z^+=G(Az^+)\le0$.
Thus $z\le0$.  No test of an unbounded vector as if it belonged
to $H$ occurs in this proof.
\end{proof}

\begin{definition}[Obstacles with finite curvature and the diffuse solution class]
\label{def:r11-diffuse-obstacle-class}
An allowed obstacle is a vector $\psi\in H$ such that
$\tau:=A\psi$ is a finite signed Radon measure and $|\tau|$ has
finite energy integral.  In particular every real original compact
smooth function is an allowed obstacle: its curvature has a bounded
density with respect to a faithful reference measure, by the
original compact coefficient control.

Fix a finite signed energy-diffuse datum $\mu$.  A solution pair
$(w,\rho)$ consists of a finite positive energy-diffuse Radon measure
$\rho$ and a signed renormalized zero-boundary solution $w$ of
\[
                               Aw=\mu+\rho
 \tag{OB5}
\]
such that
\[
                     w\ge\psi,\qquad
                     \rho\{\widetilde w>\widetilde\psi\}=0.
 \tag{OB6}
\]
The renormalized condition in (OB5) means the intrinsic class from
Definition~\ref{def:r11-signed-renormalized-boundary}, with sign
reversed: $w\in L^1(\lambda)$ is finite and quasicontinuous,
$T_kw\in H$ for every $k>0$, the original chain identity is
\[
 D(\Phi(w)S)=-\Phi'(\widetilde w)(\mu+\rho)
                         +\Phi''(\widetilde w)\Gamma_S(w)
 \tag{OB7}
\]
for $C^2$ scalar $\Phi$ with compactly supported first derivative,
and $k^{-1}\mathcal E(T_{2k}w-T_kw)\to0$.
Thus the definition does not prescribe $w$ by a Green formula.
\end{definition}

\begin{theorem}[The full finite-diffuse obstacle and contact theorem]
\label{thm:r11-diffuse-obstacle-contact}
For every allowed obstacle $\psi$ and every finite signed
energy-diffuse datum $\mu$, there is exactly one solution pair in
Definition~\ref{def:r11-diffuse-obstacle-class}.  It satisfies
\[
 \begin{split}
 w&=G(\mu+\rho),\qquad Aw=\mu+\rho,\\
 0\le\rho&\le(\tau-\mu)^+\le\tau^++\mu^-,\qquad \tau=A\psi.
 \end{split}
 \tag{OB8}
\]
The second line is the actual measure inequality, including the
contact set and every original test; no density of the arbitrary
background is assumed.  In particular the reaction is finite.

The solution is the least member of the following explicitly
nonempty class:
\[
 \mathscr M_{\mu,\psi}=
 \{G(\mu+\eta):\ \eta\ge0\text{ finite and energy diffuse},\
                         G(\mu+\eta)\ge\psi\}.
 \tag{OB9}
\]
It is therefore an actual relative envelope, with a finite positive
reaction carried by contact.  All equations are original current
equalities on the ordinary domain, for every compact smooth scalar
test and every actual moment refinement.
\end{theorem}

\noindent\textbf{Status: \texttt{NEW\_PROVED}}.  This extends the
finite-energy Hilbert obstacle problem to every finite signed diffuse
datum.  The obstacle hypothesis and part-form boundary class are
retained explicitly; it is not an unrestricted classical obstacle or
nonlinear Monge--Ampere existence assertion.

\begin{proof}
First consider a finite-energy datum $\mu_j$ and the Hilbert obstacle
solution $w_j$ from Theorem~\ref{thm:r11-relative-energy-obstacle}.
Its reaction $\rho_j$ is positive, has finite energy integral, and is
carried by $\{w_j=\psi\}$.  It need not initially be assumed finite
in total mass.  Put $f_j=w_j-\psi\in H_+$.  Its curvature is
\[
                           D(f_jS)=\tau-\mu_j-\rho_j.
\]
The total variation on the right has finite energy integral, so the
$H$ level theorem (RT2) applies even if its total mass near the
boundary is infinite.  Since $f_j^+=f_j$, its zero-level formula is
\[
        \mathbf1_{\{\widetilde f_j=0\}}
                         (\tau-\mu_j-\rho_j)\ge0.
 \tag{OB10}
\]
Combine this with the contact carrier of $\rho_j$ to obtain
\[
                          0\le\rho_j\le(\tau-\mu_j)^+.
 \tag{OB11}
\]
This is the required finite-mass bound; it is derived before being
used in any finite-data approximation argument.

Now let $\mu_j=\mathbf1_{F_j}\mu$ be an increasing finite-energy
nest for $|\mu|$, and solve these finite-energy obstacle problems.
The signed measures converge in total variation to $\mu$ and
$\mu_j^-\le\mu^-$.  Thus all reactions are dominated by the same
finite diffuse measure
\[
                         \pi=\tau^++\mu^-,\qquad 0\le\rho_j\le\pi.
 \tag{OB12}
\]
For two indices put $d=w_j-w_l\in H$.  The source equation extends
to $T_kd\in H$.  On $\{d>0\}$ the measure $\rho_j$ is zero,
because $w_j>w_l\ge\psi$ there.  On $\{d<0\}$ the measure
$\rho_l$ is zero.  Therefore
\[
 \begin{split}
 \mathcal E(T_kd)
   &=\int T_k\widetilde d\,d(\mu_j-\mu_l+\rho_j-\rho_l)\\
   &\le k\|\mu_j-\mu_l\|_{\rm TV},\qquad
 c_V\{|\widetilde d|>k\}
                         \le\frac{\|\mu_j-\mu_l\|_{\rm TV}}{k}.
 \end{split}
 \tag{OB13}
\]
The first equality follows from the closed clipping rule and the
$H$ source identity.  The reaction integral is nonpositive by the
two contact observations.  The capacity estimate follows by the
normal clipped capacity test.  Thus $w_j$ is Cauchy in capacity.

Since $w_j=G(\mu_j+\rho_j)$ by the finite-energy theorem, positivity
of the transposition Green operator and (OB12) give
\[
                  |w_j|\le G(2|\mu|+\tau^+)=:B
                                        \quad\hbox{quasi-everywhere}.
\]
The fixed potential $B$ is finite quasi-everywhere and belongs to
$L^1(\lambda)$.  The usual summable-capacity subsequence construction
from (MD7)--(MD10) gives a finite quasicontinuous limit $w$ and
capacity convergence of the whole sequence.  Coercivity implies
convergence in $\lambda$-measure; the displayed domination then
gives $w_j\to w$ in $L^1(\lambda)$.  The inequalities $w_j\ge\psi$
pass outside a polar set to $w\ge\psi$.
For completeness, the source identity also gives
$\mathcal E(T_kw_j)\le k(\|\mu\|_{\rm TV}+\pi(F))$.
Thus each sequence of bounded clippings has a weak $H$ limit, whose
$L^2$ representative is $T_kw$.  Its canonical representative is the
capacity limit: choose a summable-capacity subsequence, then apply
Mazur's lemma to its tails.  The convex combinations still converge
uniformly outside the same arbitrarily small-capacity exceptional
sets, while their strong $H$ convergence identifies the canonical
limit.  Hence $T_kw\in H$ with the same quasi-everywhere values.

Write $\rho_j=r_j\pi$ with $0\le r_j\le1$.  The finite Radon
space here is second countable, so $L^1(\pi)$ is separable.
Weak-star sequential compactness gives a subsequence with
$r_j\rightharpoonup^*r$ in $L^\infty(\pi)$, where $0\le r\le1$.
Set $\rho=r\pi$.  It is finite, positive and diffuse.  For every
$f\in L^\infty(\lambda)$ the canonical $Gf$ is bounded and hence
belongs to $L^1(\pi)$.  Passing to the limit in the finite-energy
Green identity gives
\[
 \int wf\,d\lambda
       =\lim_j\int\widetilde{Gf}\,d(\mu_j+\rho_j)
       =\int\widetilde{Gf}\,d(\mu+\rho).
 \tag{OB14}
\]
The first passage uses $L^1(\lambda)$ convergence; the two datum
passages use total variation and weak-star convergence, respectively.
Thus $w=G(\mu+\rho)$.  The signed diffuse theorem now supplies its
complete intrinsic renormalized class and every original current
equation, so no current-only uniqueness assumption is needed.
The preceding clipped representative identification makes this an
equality of canonical functions quasi-everywhere as well as in $L^1$.

For contact let $g(s)=\min(1,s^+)$.  The functions
$g(w_j-\psi)$ converge in capacity to $g(w-\psi)$ and are bounded.
Finite diffuseness of $\pi$ gives convergence in $L^1(\pi)$.
The exact finite-energy contact condition and (OB12) therefore imply
\[
 \begin{split}
 \int g(\widetilde w-\widetilde\psi)\,d\rho
 &=\lim_j\int g(\widetilde w-\widetilde\psi)\,d\rho_j\\
 &=\lim_j\int g(\widetilde w_j-\widetilde\psi)\,d\rho_j=0.
 \end{split}
 \tag{OB15}
\]
The first passage is weak-star convergence against the fixed bounded
function.  The second is bounded by the $L^1(\pi)$ error.  Since
$w\ge\psi$, this proves the exact contact carrier.  In particular
the proof does not restrict an uncontrolled weak measure limit to a
discontinuous contact set.

For any solution pair in the stated class, signed renormalized
uniqueness first gives $w=G(\mu+\rho)$.  Since $\psi=G\tau$,
the difference $f=w-\psi=G(\mu+\rho-\tau)$ is a signed diffuse
Green potential and is nonnegative.  Its signed zero-level Kato
formula gives
\[
                  \mathbf1_{\{\widetilde f=0\}}
                              (\tau-\mu-\rho)\ge0.
\]
The contact carrier yields the second line of (OB8), exactly as in
(OB10)--(OB11).

Finally let $v=G(\mu+\eta)\in\mathscr M_{\mu,\psi}$.
Then $z=w-v=G(\rho-\eta)$.  On $\{z>0\}$ one has
$w>v\ge\psi$, so $\rho$ is zero there and
$\mathbf1_{\{z>0\}}(\rho-\eta)\le0$.
Lemma~\ref{lem:r11-diffuse-positive-part-Green} gives $z\le0$.
This proves leastness.  The class is nonempty before solving the
obstacle: take $\eta=\tau^++\mu^-$, since
$G(\mu+\eta)=G(\mu^++\tau^+)\ge G\tau=\psi$.
Two solution pairs are each least, so their potentials coincide;
their original equations then identify their reaction measures.
This proves uniqueness and independence of the extracted subsequence.
\end{proof}

\begin{theorem}[Order and quantitative dependence on the signed source]
\label{thm:r11-diffuse-obstacle-stability}
For a fixed allowed obstacle, write $w_\mu$ for the solution.
For any two finite signed diffuse data $\mu,\nu$ one has
\[
 \begin{split}
 |w_\mu-w_\nu|&\le G|\mu-\nu|,\\
 \|w_\mu-w_\nu\|_{L^1(\lambda)}&\le2\|\mu-\nu\|_{\rm TV},\\
 c_V\{|\widetilde w_\mu-\widetilde w_\nu|>k\}
                              &\le\frac{\|\mu-\nu\|_{\rm TV}}{k}.
 \end{split}
 \tag{OB16}
\]
If $\mu\le\nu$, then $w_\mu\le w_\nu$.
For fixed source the solution also increases when the allowed
obstacle increases.  These are relative energy orders; no pointwise
statement outside the coefficient support is asserted.
\end{theorem}

\noindent\textbf{Status: \texttt{NEW\_PROVED}} for the finite
signed diffuse obstacle map.

\begin{proof}
Put $d=w_\mu-w_\nu$ and $b=G|\mu-\nu|\ge0$.
The signed Green potential $z=d-b$ has source
\[
            Az=\mu-\nu-|\mu-\nu|+\rho_\mu-\rho_\nu.
\]
On $\{z>0\}$ one has $d>b\ge0$, so the contact condition gives
$\rho_\mu=0$ there.  The remaining restricted source is
nonpositive.  Apply (OB3) to obtain $d\le b$; interchange the two
data to obtain $-d\le b$.  This is the first line of (OB16).
The torsion bound and the finite-datum capacity bound for $b$
give the remaining two lines.  If $\mu\le\nu$, apply (OB3)
directly to $d$, using its same contact sign, to get order.
For ordered obstacles, the solution of the larger obstacle belongs
to the class (OB9) for the smaller one, so leastness proves order.
\end{proof}

\begin{corollary}[Positive nest data give decreasing finite reactions]
\label{cor:r11-positive-obstacle-reaction-TV}
If $\mu\ge0$ and $\mu_j\uparrow\mu$ are its finite-energy nest
restrictions, then their finite-energy obstacle solutions satisfy
\[
       w_j\uparrow w_\mu,\qquad
       \rho_j\downarrow\rho_\mu,\qquad
       \|\rho_j-\rho_\mu\|_{\rm TV}\longrightarrow0,\qquad
       0\le\rho_j\le\tau^+.
 \tag{OB17}
\]
In this positive-data case the limiting reaction itself has finite
energy integral.  The potential $w_\mu$ need not belong to $H$.
\end{corollary}

\noindent\textbf{Status: \texttt{NEW\_PROVED}} as a stronger
reaction convergence statement for positive nested data.

\begin{proof}
Order gives $w_l\ge w_j$ for $l\ge j$.  On the contact set
$\{w_l=\psi\}$ one also has $w_j=\psi$.
The nonnegative $H$ vector $f=w_l-w_j$ has curvature
\[
                         D(fS)=-(\mu_l-\mu_j)-(\rho_l-\rho_j).
\]
Its zero-level Kato inequality, restricted to the contact set of
$w_l$, gives $\rho_l\le\rho_j$ there because $\mu_l\ge\mu_j$.
Off that set $\rho_l=0$.  Hence the reactions decrease as measures.
Their common finite bound $\tau^+$ from (OB11) gives total-variation
convergence.  The preceding uniqueness and source convergence identify
the limits with $w_\mu,\rho_\mu$.  Domination by the finite-energy
measure $\tau^+$ gives finite energy integral for the reaction.
This does not put the possibly infinite-energy source $\mu$ or the
whole potential $w_\mu$ in $H$.
\end{proof}

\paragraph{Dependency and scope audit.}
The finite-energy obstacle/contact theorem is used only for the nest
approximants.  Its level-zero Kato consequence proves the uniform
finite reaction bound before extracting measures.  The signed
renormalized theorem supplies both the limiting boundary class and
the positive-part Green representation used in comparison; the
unbounded potentials are never tested as if they were $H$ vectors.
The common finite diffuse domination is what passes contact through
the signed approximation.  No ambient regularization of an arbitrary
upper envelope, no current pullback under an auxiliary field, and no
analytic PS/PL quotient is used.  Removing the finite-curvature
obstacle hypothesis, imposing ordinary pointwise boundary values, or
realizing these envelopes by ambient classical psh functions requires
additional arguments and is not asserted here.
